\documentclass[12pt,a4paper]{article}
\usepackage[utf8]{inputenc}
\usepackage[T1]{fontenc}
\usepackage[spanish]{babel}
\usepackage{geometry}
\usepackage{booktabs}
\usepackage{longtable}
\usepackage{array}
\usepackage{xcolor}
\usepackage{listings}
\usepackage{hyperref}
\usepackage{fancyhdr}
\usepackage{graphicx}
\usepackage{titlesec}
\usepackage{enumitem}
\usepackage{mdframed}
\usepackage{tabularx}
\usepackage{multirow}
\usepackage{amsmath}
\usepackage{amssymb}

\geometry{margin=2.5cm}

% ---- Colores del sistema ----
\definecolor{primaryblue}{RGB}{33,150,243}
\definecolor{criticalred}{RGB}{244,67,54}
\definecolor{normalgreen}{RGB}{76,175,80}
\definecolor{warningorange}{RGB}{255,152,0}
\definecolor{codegray}{RGB}{245,245,245}
\definecolor{darkgray}{RGB}{55,55,55}

\lstset{
  backgroundcolor=\color{codegray},
  basicstyle=\ttfamily\small,
  breaklines=true,
  frame=single,
  rulecolor=\color{gray},
  keywordstyle=\color{primaryblue}\bfseries,
  commentstyle=\color{normalgreen},
  stringstyle=\color{criticalred},
  showstringspaces=false,
  tabsize=2,
  captionpos=b,
  numbers=left,
  numberstyle=\tiny\color{gray},
}

\pagestyle{fancy}
\fancyhf{}
\rhead{EdgeLeak AI -- Reporte Técnico Completo}
\lhead{Documentación Integral}
\cfoot{\thepage}

\hypersetup{
  colorlinks=true,
  linkcolor=primaryblue,
  urlcolor=primaryblue,
  citecolor=primaryblue,
  pdftitle={Reporte Técnico Completo - EdgeLeak AI},
  pdfauthor={Proyecto EdgeLeak AI},
}

\titleformat{\section}{\Large\bfseries\color{primaryblue}}{}{0em}{\thesection.\enspace}
\titleformat{\subsection}{\large\bfseries\color{darkgray}}{}{0em}{\thesubsection.\enspace}
\titleformat{\subsubsection}{\normalsize\bfseries}{}{0em}{\thesubsubsection.\enspace}

% ====================================================================
\title{
  \Huge\textbf{Reporte Técnico Completo}\\[0.5em]
  \textcolor{primaryblue}{\rule{\linewidth}{1.5pt}}\\[0.3em]
  \Large\textbf{EdgeLeak AI}\\[0.2em]
  \large Sistema Inteligente de Detección de Fugas de Agua\\[0.5em]
  \normalsize Compilación integral de todas las fases de documentación técnica:\\
  Fase MVP $\rightarrow$ Fase 2 $\rightarrow$ Fase 3 (CRUD/Roles) $\rightarrow$
  Fase 4 (Correcciones) $\rightarrow$ Actualización \#5
}
\author{Proyecto EdgeLeak AI}
\date{\today}
% ====================================================================

\begin{document}

\maketitle
\thispagestyle{fancy}

\begin{mdframed}[linecolor=primaryblue,linewidth=1.5pt,backgroundcolor=blue!3]
\textbf{Alcance del documento:} Este reporte consolida \emph{toda} la documentación técnica
del proyecto EdgeLeak AI desde su Fase MVP hasta la Actualización~\#5. Incluye la
arquitectura del sistema, el stack tecnológico, los esquemas de base de datos, los módulos
de seguridad, la integración con IA, los flujos de autenticación, el módulo CRUD
administrativo, las pruebas unitarias, y las correcciones implementadas en cada iteración.
Todos los datos están organizados en tablas para facilitar la auditoría y el mantenimiento.
\end{mdframed}

\vspace{1em}
\tableofcontents
\newpage

% ====================================================================
\section{Visión General del Proyecto}
% ====================================================================

EdgeLeak AI es una aplicación móvil (Android/iOS) desarrollada en Flutter que implementa
un sistema de detección inteligente de fugas de agua. Ante la ausencia de hardware físico
en las fases iniciales, el sistema incorpora un gemelo digital (simulador IoT) y delega
el análisis de patrones a un modelo de lenguaje de gran escala (LLM) a través de la API
de Groq.

\begin{longtable}{>{\bfseries}p{4cm} p{9.5cm}}
\caption{Ficha técnica del proyecto EdgeLeak AI} \label{tab:ficha-tecnica} \\
\toprule
\textbf{Atributo} & \textbf{Valor} \\
\midrule
\endfirsthead
\toprule
\textbf{Atributo} & \textbf{Valor} \\
\midrule
\endhead
\bottomrule
\endlastfoot
Nombre & EdgeLeak AI \\
Plataformas objetivo & Android, iOS (con soporte potencial a Web/Desktop vía Flutter) \\
Framework & Flutter SDK / Lenguaje Dart \\
IA & Groq API, Modelo Llama-3.3-70B-Versatile \\
Persistencia & SQLite local (paquete \texttt{sqflite}) \\
Correo & Google Workspace SMTP (SSL/TLS) \\
Seguridad de entorno & \texttt{flutter\_dotenv} (variables de entorno) \\
Patrón de arquitectura & MVC adaptado con \texttt{ChangeNotifier} / \texttt{ListenableBuilder} \\
Metodología de pruebas & TDD (Test-Driven Development) con \texttt{flutter\_test} \\
Versión actual de BD & \texttt{edgeleak\_v3.db} (schema version 3) \\
\end{longtable}

% ====================================================================
\section{Arquitectura del Sistema (MVC Reactivo)}
% ====================================================================

El proyecto sigue el patrón MVC adaptado a las convenciones de reactividad de Flutter,
garantizando bajo acoplamiento y alta cohesión.

\begin{longtable}{>{\ttfamily}p{3.5cm} p{3.5cm} p{7cm}}
\caption{Estructura de capas del proyecto (directorio \texttt{lib/})} \label{tab:arquitectura} \\
\toprule
\textbf{Directorio} & \textbf{Capa MVC} & \textbf{Responsabilidad} \\
\midrule
\endfirsthead
\toprule
\textbf{Directorio} & \textbf{Capa MVC} & \textbf{Responsabilidad} \\
\midrule
\endhead
\bottomrule
\endlastfoot
/config          & Configuración & Enrutamiento estático (\texttt{AppRoutes}), paleta de colores semánticos (Material Design 3), temas globales. \\[0.4em]
/data/models     & Modelo        & POJOs en Dart: \texttt{UsuarioModel}, \texttt{AlertaFugaModel}, \texttt{LecturaSensorModel}. Incluyen serialización/deserialización JSON/Map. \\[0.4em]
/data/services   & Modelo        & Capa de infraestructura: \texttt{DatabaseService} (SQLite), \texttt{EmailService} (SMTP), \texttt{GroqApiService} (REST). \\[0.4em]
/controllers     & Controlador   & \texttt{AuthController} (sesión, registro, recuperación), \texttt{DashboardController} (simulador IoT, integración IA). \\[0.4em]
/view/screens    & Vista         & Pantallas: Login, Registro, Dashboard, Historial, Recuperación, Cambio de Clave, Administración. \\[0.4em]
/view/widgets    & Vista         & Componentes reutilizables: \texttt{StatusIndicatorWidget}, \texttt{WaterWaveWidget}. \\
\end{longtable}

\subsection{Enrutamiento Estático (\texttt{AppRoutes})}

\begin{longtable}{>{\ttfamily}p{4.5cm} >{\ttfamily}p{4.5cm} p{4.5cm}}
\caption{Rutas definidas en \texttt{AppRoutes}} \label{tab:rutas} \\
\toprule
\textbf{Constante} & \textbf{Valor (String)} & \textbf{Pantalla} \\
\midrule
\endfirsthead
\toprule
\textbf{Constante} & \textbf{Valor (String)} & \textbf{Pantalla} \\
\midrule
\endhead
\bottomrule
\endlastfoot
loginScreen          & 'login'           & \texttt{LoginScreen} \\
registerScreen       & 'register'        & \texttt{RegisterScreen} \\
forgotPasswordScreen & 'forgot\_password' & \texttt{ForgotPasswordScreen} \\
changePasswordScreen & 'change\_password' & \texttt{ChangePasswordScreen} \\
dashboardScreen      & 'dashboard'       & \texttt{DashboardScreen} \\
historialScreen      & 'historial'       & \texttt{HistorialScreen} \\
adminUsersScreen     & 'admin\_users'     & \texttt{AdminUsersScreen} \\
\end{longtable}

% ====================================================================
\section{Stack Tecnológico}
% ====================================================================

\begin{longtable}{p{3.5cm} p{3.5cm} p{7cm}}
\caption{Stack tecnológico completo del sistema} \label{tab:stack} \\
\toprule
\textbf{Categoría} & \textbf{Tecnología} & \textbf{Uso y Justificación} \\
\midrule
\endfirsthead
\toprule
\textbf{Categoría} & \textbf{Tecnología} & \textbf{Uso y Justificación} \\
\midrule
\endhead
\bottomrule
\endlastfoot
Frontend & Flutter SDK & Framework multiplataforma. Renderizado nativo de alta performance. \\[0.3em]
Lenguaje & Dart & Tipado estático, null-safety, async/await, ecosistema pub.dev. \\[0.3em]
IA/LLM & Groq API & Inferencia ultrarrápida con Llama-3.3-70B. Temperatura 0 para determinismo. \\[0.3em]
Persistencia & sqflite + SQLite & Base de datos relacional embebida. Funciona en modo offline. \\[0.3em]
Correo & mailer + Gmail SMTP & Envío de contraseñas temporales por SSL/TLS. \\[0.3em]
Variables de entorno & flutter\_dotenv & Separación de credenciales del código fuente. \\[0.3em]
Testing & flutter\_test & Pruebas unitarias y de widget bajo metodología TDD. \\[0.3em]
Rutas de archivos & path (paquete) & Construcción multiplataforma de rutas de archivos SQLite. \\
\end{longtable}

% ====================================================================
\section{Esquema de Base de Datos (Evolución Histórica)}
% ====================================================================

La base de datos local ha evolucionado a través de múltiples versiones de schema para
soportar nuevas funcionalidades.

\subsection{Historial de Versiones del Schema}

\begin{longtable}{p{2cm} p{3.5cm} p{8cm}}
\caption{Historial de migraciones de la base de datos} \label{tab:historial-db} \\
\toprule
\textbf{Versión} & \textbf{Archivo DB} & \textbf{Cambios Introducidos} \\
\midrule
\endfirsthead
\toprule
\textbf{Versión} & \textbf{Archivo DB} & \textbf{Cambios Introducidos} \\
\midrule
\endhead
\bottomrule
\endlastfoot
v1 & edgeleak\_v2.db & Schema inicial: tablas \texttt{historial} y \texttt{usuarios} con campos \texttt{id}, \texttt{nombre}, \texttt{correo}, \texttt{password}. \\[0.3em]
v2 & edgeleak\_v3.db & Adición de columna \texttt{es\_temporal} (INTEGER 0/1) en \texttt{usuarios} para soporte de contraseñas temporales. Adición de columna \texttt{rol} (TEXT) para RBAC. \\[0.3em]
v3 & edgeleak\_v3.db & Adición de columnas \texttt{primer\_nombre}, \texttt{segundo\_nombre}, \texttt{primer\_apellido}, \texttt{segundo\_apellido}. Migración automática con preservación de datos legados. \\
\end{longtable}

\subsection{Tabla: \texttt{usuarios} (Schema v3 actual)}

\begin{longtable}{>{\ttfamily}p{3.8cm} p{2.5cm} p{2cm} p{5.2cm}}
\caption{Definición completa de la tabla \texttt{usuarios}} \label{tab:tabla-usuarios} \\
\toprule
\textbf{Columna} & \textbf{Tipo SQL} & \textbf{Restricción} & \textbf{Descripción} \\
\midrule
\endfirsthead
\toprule
\textbf{Columna} & \textbf{Tipo SQL} & \textbf{Restricción} & \textbf{Descripción} \\
\midrule
\endhead
\bottomrule
\endlastfoot
id              & INTEGER & PK, AUTOINCR. & Identificador único autoincremental. \\[0.3em]
primer\_nombre  & TEXT    & DEFAULT ''    & Primer nombre del usuario (obligatorio en UI). \\[0.3em]
segundo\_nombre & TEXT    & DEFAULT ''    & Segundo nombre (opcional). \\[0.3em]
primer\_apellido& TEXT    & DEFAULT ''    & Primer apellido (obligatorio en UI). \\[0.3em]
segundo\_apellido& TEXT   & DEFAULT ''    & Segundo apellido (opcional). \\[0.3em]
correo          & TEXT    & UNIQUE        & Correo electrónico. Llave de negocio para evitar duplicidad. \\[0.3em]
password        & TEXT    & ---           & Contraseña en texto plano (se planea migrar a bcrypt). \\[0.3em]
es\_temporal    & INTEGER & DEFAULT 0     & Flag de contraseña temporal. 0=permanente, 1=temporal. \\[0.3em]
rol             & TEXT    & DEFAULT 'operador' & Privilegio del usuario: \texttt{'admin'} o \texttt{'operador'}. \\
\end{longtable}

\subsection{Tabla: \texttt{historial}}

\begin{longtable}{>{\ttfamily}p{3cm} p{3cm} p{7.5cm}}
\caption{Definición de la tabla \texttt{historial}} \label{tab:tabla-historial} \\
\toprule
\textbf{Columna} & \textbf{Tipo SQL} & \textbf{Descripción} \\
\midrule
\endfirsthead
\toprule
\textbf{Columna} & \textbf{Tipo SQL} & \textbf{Descripción} \\
\midrule
\endhead
\bottomrule
\endlastfoot
id        & INTEGER & PK, AUTOINCREMENTAL. \\
veredicto & TEXT    & Clasificación de la IA (Ej. ``Fuga Detectada'', ``Normal''). \\
severidad & TEXT    & Nivel de criticidad: \texttt{"Normal"}, \texttt{"Advertencia"}, \texttt{"Crítica"}. \\
mensaje   & TEXT    & Acción recomendada generada por el LLM. \\
fecha     & TEXT    & Timestamp en formato ISO-8601. \\
\end{longtable}

% ====================================================================
\section{Modelo de Datos: \texttt{UsuarioModel}}
% ====================================================================

\begin{longtable}{>{\ttfamily}p{4cm} p{2.5cm} p{3cm} p{4cm}}
\caption{Propiedades del modelo \texttt{UsuarioModel}} \label{tab:usuario-model} \\
\toprule
\textbf{Propiedad Dart} & \textbf{Tipo} & \textbf{Requerido} & \textbf{Descripción} \\
\midrule
\endfirsthead
\toprule
\textbf{Propiedad Dart} & \textbf{Tipo} & \textbf{Requerido} & \textbf{Descripción} \\
\midrule
\endhead
\bottomrule
\endlastfoot
id              & int?   & No  & Clave primaria (null para nuevos registros). \\
primerNombre    & String & Sí  & Primer nombre. \\
segundoNombre   & String & No  & Segundo nombre (def. \texttt{''}). \\
primerApellido  & String & Sí  & Primer apellido. \\
segundoApellido & String & No  & Segundo apellido (def. \texttt{''}). \\
correo          & String & Sí  & Email validado con regex en la UI. \\
password        & String & Sí  & Contraseña. \\
esTemporal      & int    & No  & \texttt{0}=permanente, \texttt{1}=temporal (def. \texttt{0}). \\
rol             & String & No  & \texttt{'admin'} o \texttt{'operador'} (def. \texttt{'operador'}). \\
\midrule
\multicolumn{4}{l}{\textbf{Getters calculados:}} \\
\midrule
nombre  & String & --- & Nombre completo calculado a partir de los 4 campos. \\
esAdmin & bool   & --- & \texttt{true} si \texttt{rol == 'admin'}. \\
\end{longtable}

% ====================================================================
\section{Módulo de Autenticación (\texttt{AuthController})}
% ====================================================================

\begin{longtable}{>{\ttfamily}p{5.5cm} p{8cm}}
\caption{Métodos públicos de \texttt{AuthController}} \label{tab:auth-controller} \\
\toprule
\textbf{Método / Getter} & \textbf{Descripción} \\
\midrule
\endfirsthead
\toprule
\textbf{Método / Getter} & \textbf{Descripción} \\
\midrule
\endhead
\bottomrule
\endlastfoot
\texttt{login(correo, pass)} &
  Autentica al usuario contra SQLite. Si \texttt{es\_temporal == 1}, la UI redirige a \texttt{ChangePasswordScreen}. \\[0.3em]
\texttt{register(p1, p2, ap1, ap2, correo, pass, \{rol\})} &
  Registra un nuevo usuario con los 4 campos de nombre. Valida política de contraseña antes de insertar. \\[0.3em]
\texttt{enviarCorreoRecuperacion(correo)} &
  Genera contraseña temporal, la persiste con \texttt{es\_temporal=1} y la envía vía SMTP. \\[0.3em]
\texttt{cambiarPasswordObligatorio(pass)} &
  Actualiza la contraseña en BD, \textbf{limpia} \texttt{usuarioActual = null} y retorna \texttt{true}. \\[0.3em]
\texttt{obtenerTodosLosUsuarios()} &
  Carga la lista completa de usuarios en \texttt{\_listaUsuarios} (para CRUD Admin). \\[0.3em]
\texttt{eliminarUsuario(id)} &
  Borra un usuario. Impide autoeliminación validando contra \texttt{usuarioActual?.id}. \\[0.3em]
\texttt{actualizarDatosUsuario(usuario)} &
  Actualiza nombre y rol de un usuario existente. \\[0.3em]
\texttt{validatePassword(pass)} &
  Actualiza los 4 flags de validación de contraseña en tiempo real. \\[0.3em]
\texttt{generarContrasenaSegura()} &
  Genera una contraseña de 12 caracteres con shuffle criptográfico. \\[0.3em]
\texttt{isPasswordValid} &
  Getter: \texttt{true} si los 4 criterios de seguridad se cumplen. \\[0.3em]
\texttt{resetValidators()} &
  Reinicia los validadores de contraseña y el mensaje de error. \\
\end{longtable}

\subsection{Criterios de Validación de Contraseña}

\begin{longtable}{p{3cm} p{5cm} p{5.5cm}}
\caption{Criterios de validación de contraseña fuerte} \label{tab:criterios-pass} \\
\toprule
\textbf{Criterio} & \textbf{Regex / Condición} & \textbf{Descripción} \\
\midrule
\endfirsthead
\toprule
\textbf{Criterio} & \textbf{Regex / Condición} & \textbf{Descripción} \\
\midrule
\endhead
\bottomrule
\endlastfoot
Longitud mínima   & \texttt{length >= 8}            & Mínimo 8 caracteres. \\
Mayúsculas        & \texttt{RegExp(r'[A-Z]')}        & Al menos 1 letra mayúscula. \\
Números           & \texttt{RegExp(r'[0-9]')}        & Al menos 1 dígito numérico. \\
Caracteres especiales & \texttt{RegExp(r'[!@\#\$\&*~]')} & Al menos 1 símbolo del conjunto \texttt{!@\#\$\&*\textasciitilde}. \\
\end{longtable}

% ====================================================================
\section{Módulo de Simulación IoT (Gemelo Digital)}
% ====================================================================

El \texttt{DashboardController} implementa un simulador de sensor de caudal (YF-S201)
mediante funciones de \texttt{dart:math}, generando lecturas periódicas cada 2 segundos.

\begin{longtable}{p{3.5cm} p{3.5cm} p{3.5cm} p{3cm}}
\caption{Estados del simulador de caudal (Gemelo Digital)} \label{tab:simulador} \\
\toprule
\textbf{Modo} & \textbf{Rango de Caudal} & \textbf{Representación Visual} & \textbf{Veredicto IA} \\
\midrule
\endfirsthead
\toprule
\textbf{Modo} & \textbf{Rango de Caudal} & \textbf{Representación Visual} & \textbf{Veredicto IA} \\
\midrule
\endhead
\bottomrule
\endlastfoot
Normal    & 0.1 -- 0.5 L/min & Ondas lentas, baja amplitud  & Normal \\
Anomalía  & 1.2 -- 2.5 L/min & Ondas medias, amp. moderada  & Advertencia \\
Fuga Crítica & 6.0 -- 9.0 L/min & Ondas rápidas, alta amplitud & Crítico \\
\end{longtable}

Las fluctuaciones se visualizan mediante un \texttt{CustomPainter} que renderiza ondas
senoidales: $y = A \cdot \sin(\omega x + \phi)$, donde la amplitud $A$ y la velocidad
angular $\omega$ varían según el estado del sistema.

% ====================================================================
\section{Integración de Inteligencia Artificial (Groq API)}
% ====================================================================

\begin{longtable}{p{4cm} p{9.5cm}}
\caption{Configuración del payload enviado a Groq API} \label{tab:groq-payload} \\
\toprule
\textbf{Parámetro} & \textbf{Valor / Descripción} \\
\midrule
\endfirsthead
\toprule
\textbf{Parámetro} & \textbf{Valor / Descripción} \\
\midrule
\endhead
\bottomrule
\endlastfoot
Modelo & \texttt{llama-3.3-70b-versatile} \\
Método HTTP & \texttt{POST} a \texttt{https://api.groq.com/openai/v1/chat/completions} \\
Response format & \texttt{json\_object} (garantiza respuesta parseable) \\
Temperature & \texttt{0.0} (determinismo: elimina alucinaciones y variabilidad) \\
System Role & Define esquema JSON exacto esperado: \texttt{\{veredicto, severidad, mensaje\}} \\
User Role & Inyecta lectura de caudal en L/min y parámetros del sistema métrico \\
\end{longtable}

\begin{longtable}{p{3cm} p{3cm} p{7.5cm}}
\caption{Estructura del objeto JSON retornado por la IA} \label{tab:json-ia} \\
\toprule
\textbf{Campo} & \textbf{Tipo} & \textbf{Descripción} \\
\midrule
\endfirsthead
\toprule
\textbf{Campo} & \textbf{Tipo} & \textbf{Descripción} \\
\midrule
\endhead
\bottomrule
\endlastfoot
\texttt{veredicto} & String & Clasificación: \texttt{"Normal"}, \texttt{"Goteo/Anomalía"} o \texttt{"Fuga Detectada"}. \\
\texttt{severidad} & String & Nivel: \texttt{"Normal"}, \texttt{"Advertencia"} o \texttt{"Crítica"}. \\
\texttt{mensaje}   & String & Recomendación accionable generada por el LLM en español. \\
\end{longtable}

% ====================================================================
\section{Protocolo de Recuperación de Credenciales (SMTP)}
% ====================================================================

\begin{longtable}{p{4cm} p{9.5cm}}
\caption{Características del módulo de recuperación SMTP} \label{tab:smtp} \\
\toprule
\textbf{Característica} & \textbf{Descripción Técnica} \\
\midrule
\endfirsthead
\toprule
\textbf{Característica} & \textbf{Descripción Técnica} \\
\midrule
\endhead
\bottomrule
\endlastfoot
Protocolo & SMTP (Simple Mail Transfer Protocol) directo a Google Workspace. \\
Seguridad de transporte & SSL/TLS implícito (cifrado extremo a extremo). \\
Autenticación S2S & Contraseñas de Aplicación de Google con 2FA activo. Credenciales aisladas del código fuente vía \texttt{.env}. \\
Formato de correo & Plantilla HTML/CSS responsiva (no texto plano). \\
Cuenta remitente & \texttt{edgeleak.ai@gmail.com} (configurada en variables de entorno). \\
Algoritmo de generación & \texttt{generarContrasenaSegura()}: longitud 12, inyección obligatoria de 4 tipos de caracteres + dispersión aleatoria. \\
\end{longtable}

\subsection{Algoritmo de Generación de Contraseñas Seguras}

El método \texttt{generarContrasenaSegura()} opera en tres fases:

\begin{enumerate}[leftmargin=2em]
  \item \textbf{Inyección obligatoria:} Garantiza al menos un carácter de cada conjunto:
        Mayúsculas \texttt{[A-Z]}, Minúsculas \texttt{[a-z]}, Números \texttt{[0-9]},
        Símbolos \texttt{[!@\#\$\&*\textasciitilde]}.
  \item \textbf{Relleno aleatorio:} Completa hasta longitud estándar de 12 caracteres
        usando el conjunto unificado.
  \item \textbf{Dispersión (Shuffle):} Aplica permutación aleatoria sobre los índices del
        arreglo para eliminar patrones predecibles.
\end{enumerate}

% ====================================================================
\section{Flujo de Intercepción UX: Cambio Obligatorio de Contraseña}
% ====================================================================

\begin{longtable}{p{0.5cm} p{4cm} p{4cm} p{5cm}}
\caption{Flujo completo de contraseña temporal} \label{tab:flujo-temp} \\
\toprule
\textbf{\#} & \textbf{Acción} & \textbf{Clase/Método} & \textbf{Resultado} \\
\midrule
\endfirsthead
\toprule
\textbf{\#} & \textbf{Acción} & \textbf{Clase/Método} & \textbf{Resultado} \\
\midrule
\endhead
\bottomrule
\endlastfoot
1 & Usuario solicita recuperación & \texttt{ForgotPasswordScreen} & Se busca el correo en BD. \\
2 & Sistema genera clave temporal & \texttt{generarContrasenaSegura()} & Cadena de 12 chars segura. \\
3 & Se persiste con flag temporal & \texttt{actualizarPassword(id, pass, true)} & \texttt{es\_temporal = 1} en BD. \\
4 & Se envía por correo & \texttt{EmailService.enviarCorreoRecuperacion()} & HTML/CSS vía SMTP. \\
5 & Usuario inicia sesión & \texttt{AuthController.login()} & Lee \texttt{es\_temporal}. \\
6 & Intercepción de navegación & \texttt{LoginScreen} & Redirige a \texttt{ChangePasswordScreen}. \\
7 & Pantalla bloqueada & \texttt{PopScope(canPop: false)} & No permite retroceder. \\
8 & Usuario cambia contraseña & \texttt{cambiarPasswordObligatorio()} & \texttt{es\_temporal = 0}, \texttt{usuarioActual = null}. \\
9 & Redirección segura & \texttt{ChangePasswordScreen} & \texttt{pushNamedAndRemoveUntil} → Login. \\
10 & Usuario se autentica de nuevo & \texttt{LoginScreen} & Acceso al Dashboard con credencial permanente. \\
\end{longtable}

% ====================================================================
\section{Control de Acceso Basado en Roles (RBAC)}
% ====================================================================

\begin{longtable}{p{3cm} p{5.5cm} p{5cm}}
\caption{Definición de roles y sus privilegios} \label{tab:rbac} \\
\toprule
\textbf{Rol} & \textbf{Funcionalidades Habilitadas} & \textbf{Restricciones} \\
\midrule
\endfirsthead
\toprule
\textbf{Rol} & \textbf{Funcionalidades Habilitadas} & \textbf{Restricciones} \\
\midrule
\endhead
\bottomrule
\endlastfoot
\texttt{operador} &
  Dashboard, Historial de eventos, Análisis IA, Simulador. &
  Sin acceso al CRUD de usuarios (\texttt{AdminUsersScreen} no se renderiza). \\[0.3em]
\texttt{admin} &
  Todo lo de operador + Gestión de usuarios (crear/editar/eliminar), escalado de privilegios. &
  No puede eliminar su propia cuenta activa. Sudo Mode requerido para operaciones críticas. \\
\end{longtable}

\subsection{Inyección Segura de Privilegios (Registro de Administradores)}

\begin{longtable}{p{4.5cm} p{9cm}}
\caption{Flujo de asignación de rol en el registro} \label{tab:asignacion-rol} \\
\toprule
\textbf{Condición} & \textbf{Resultado} \\
\midrule
\endfirsthead
\toprule
\textbf{Condición} & \textbf{Resultado} \\
\midrule
\endhead
\bottomrule
\endlastfoot
\texttt{SECRET\_ADMIN\_CODE} en \texttt{.env} coincide con el código ingresado en el campo opcional &
  \texttt{UsuarioModel} instanciado con \texttt{rol: 'admin'}. \\
Código no coincide o campo vacío &
  \texttt{UsuarioModel} instanciado con \texttt{rol: 'operador'} (mínima autoridad por defecto). \\
\end{longtable}

% ====================================================================
\section{Módulo CRUD Administrativo (\texttt{AdminUsersScreen})}
% ====================================================================

\begin{longtable}{p{4cm} p{9.5cm}}
\caption{Funcionalidades del panel de administración de usuarios} \label{tab:crud-admin} \\
\toprule
\textbf{Función} & \textbf{Implementación} \\
\midrule
\endfirsthead
\toprule
\textbf{Función} & \textbf{Implementación} \\
\midrule
\endhead
\bottomrule
\endlastfoot
Listado reactivo & \texttt{ListenableBuilder} sincronizado con \texttt{AuthController.listaUsuarios}. \\[0.3em]
Búsqueda dinámica & Filtrado en memoria con \texttt{.where()} aplicado sobre \texttt{listaUsuarios}. Normalización con \texttt{toLowerCase()}. \\[0.3em]
Identificación visual & Avatares \textcolor{criticalred}{\textbf{Rojo}} = admin, \textcolor{primaryblue}{\textbf{Azul}} = operador. \\[0.3em]
Edición (Sudo Mode) & AlertDialog con 4 campos de nombre + rol + contraseña de confirmación. \\[0.3em]
Eliminación (Sudo Mode) & Diálogo de confirmación con validación de contraseña del administrador activo. \\[0.3em]
Seguridad de autoeliminación & \texttt{eliminarUsuario()} compara \texttt{usuario.id} contra \texttt{usuarioActual?.id} antes de ejecutar DELETE. \\
\end{longtable}

\subsection{Defensa en Profundidad: Sudo Mode}

El mecanismo de reautenticación activa opera de la siguiente forma para proteger operaciones
críticas (editar rol, eliminar usuario):

\begin{enumerate}[leftmargin=2em]
  \item El administrador inicia la operación (botón editar o eliminar).
  \item Se abre un \texttt{AlertDialog} con \texttt{StatefulBuilder} que solicita la
        contraseña actual del administrador.
  \item El controlador compara \texttt{input == usuarioActual?.password}.
  \item \textbf{Si coincide:} se ejecuta la transacción SQL y se notifica el resultado.
  \item \textbf{Si no coincide:} se muestra un \texttt{SnackBar} de error y se cancela
        la operación sin acceder a la base de datos.
\end{enumerate}

% ====================================================================
\section{Formulario de Registro (Diseño Actual -- Actualización \#5)}
% ====================================================================

\begin{longtable}{p{4.5cm} p{2.5cm} p{3cm} p{3.5cm}}
\caption{Campos del formulario de registro actualizado} \label{tab:form-registro-v5} \\
\toprule
\textbf{Campo} & \textbf{Oblig.} & \textbf{Tipo de entrada} & \textbf{Validación} \\
\midrule
\endfirsthead
\toprule
\textbf{Campo} & \textbf{Oblig.} & \textbf{Tipo de entrada} & \textbf{Validación} \\
\midrule
\endhead
\bottomrule
\endlastfoot
Primer Nombre *        & Sí & Texto libre & No vacío \\
Segundo Nombre         & No & Texto libre & Ninguna \\
Primer Apellido *      & Sí & Texto libre & No vacío \\
Segundo Apellido       & No & Texto libre & Ninguna \\
Correo Electrónico *   & Sí & Email       & Regex RFC-5322 simplificado \\
Contraseña *           & Sí & Contraseña  & 4 criterios de seguridad \\
Código Admin           & No & Texto libre & Comparación contra \texttt{.env} \\
\end{longtable}

\subsection{Validación de Correo Electrónico}

\begin{longtable}{p{5cm} p{2.5cm} p{6cm}}
\caption{Ejemplos de validación de correo electrónico con regex} \label{tab:email-validacion-full} \\
\toprule
\textbf{Cadena de entrada} & \textbf{Resultado} & \textbf{Motivo} \\
\midrule
\endfirsthead
\toprule
\textbf{Cadena de entrada} & \textbf{Resultado} & \textbf{Motivo} \\
\midrule
\endhead
\bottomrule
\endlastfoot
\texttt{usuario@dominio.com}   & \textcolor{normalgreen}{\checkmark} Válido   & Formato completo RFC. \\
\texttt{user.name+tag@sub.d.co}& \textcolor{normalgreen}{\checkmark} Válido   & Caracteres especiales permitidos. \\
\texttt{sinArroba.com}         & \textcolor{criticalred}{\texttimes} Inválido  & Ausencia de \texttt{@}. \\
\texttt{@sinusuario.com}       & \textcolor{criticalred}{\texttimes} Inválido  & Parte local vacía. \\
\texttt{usuario@}              & \textcolor{criticalred}{\texttimes} Inválido  & Dominio ausente. \\
\texttt{usuario@dom}           & \textcolor{criticalred}{\texttimes} Inválido  & TLD < 2 caracteres. \\
\texttt{(vacío)}               & \textcolor{criticalred}{\texttimes} Inválido  & Campo obligatorio vacío. \\
\end{longtable}

% ====================================================================
\section{Flujos de Navegación del Sistema}
% ====================================================================

\begin{longtable}{p{5cm} p{4cm} p{4.5cm}}
\caption{Flujos de navegación completos del sistema} \label{tab:flujos-nav} \\
\toprule
\textbf{Acción del usuario} & \textbf{Pantalla origen} & \textbf{Destino (post-\#5)} \\
\midrule
\endfirsthead
\toprule
\textbf{Acción del usuario} & \textbf{Pantalla origen} & \textbf{Destino (post-\#5)} \\
\midrule
\endhead
\bottomrule
\endlastfoot
Iniciar aplicación & --- & \texttt{LoginScreen} (ruta inicial) \\
Login exitoso (clave normal) & \texttt{LoginScreen} & \texttt{DashboardScreen} \\
Login con clave temporal & \texttt{LoginScreen} & \texttt{ChangePasswordScreen} \\
Cambio de clave exitoso & \texttt{ChangePasswordScreen} & \texttt{LoginScreen} (stack limpio) \\
Registro exitoso & \texttt{RegisterScreen} & \texttt{LoginScreen} (stack limpio) \\
Click ``¿Olvidaste tu contraseña?'' & \texttt{LoginScreen} & \texttt{ForgotPasswordScreen} \\
Recuperación exitosa & \texttt{ForgotPasswordScreen} & \texttt{LoginScreen} \\
Acceso al historial & \texttt{DashboardScreen} & \texttt{HistorialScreen} \\
Acceso al CRUD (solo admin) & \texttt{DashboardScreen} & \texttt{AdminUsersScreen} \\
\end{longtable}

% ====================================================================
\section{Pruebas Automatizadas (TDD)}
% ====================================================================

\subsection{Estrategia de Testing}

\begin{longtable}{p{4cm} p{9.5cm}}
\caption{Enfoques de prueba implementados} \label{tab:estrategia-test} \\
\toprule
\textbf{Tipo de Prueba} & \textbf{Descripción} \\
\midrule
\endfirsthead
\toprule
\textbf{Tipo de Prueba} & \textbf{Descripción} \\
\midrule
\endhead
\bottomrule
\endlastfoot
Pruebas Unitarias & Auditan la lógica pura de los controladores (sin dependencias de Flutter). Archivo: \texttt{test/unit\_tests/user\_logic\_test.dart}. \\[0.3em]
Pruebas de Widget & Simulan entorno virtual de Flutter con \texttt{WidgetTester}. Montan pantallas y verifican el árbol de widgets. Archivo: \texttt{test/widget\_test.dart}. \\[0.3em]
Manejo de memoria & Uso de \texttt{addTearDown()} para eliminar instancias activas de \texttt{Timer.periodic}, previniendo fugas de memoria en la Dart VM. \\
\end{longtable}

\subsection{Casos de Prueba Unitaria (Actualizados a v5)}

\begin{longtable}{p{6cm} p{2.5cm} p{5cm}}
\caption{Casos de prueba en \texttt{user\_logic\_test.dart}} \label{tab:unit-tests} \\
\toprule
\textbf{Nombre del Test} & \textbf{Estado} & \textbf{Valida} \\
\midrule
\endfirsthead
\toprule
\textbf{Nombre del Test} & \textbf{Estado} & \textbf{Valida} \\
\midrule
\endhead
\bottomrule
\endlastfoot
El modelo debe detectar correctamente el rol admin &
  \textcolor{normalgreen}{\checkmark} Activo &
  \texttt{esAdmin == true} para \texttt{rol='admin'}. \\[0.3em]
El modelo debe denegar acceso admin a un operador &
  \textcolor{normalgreen}{\checkmark} Activo &
  \texttt{esAdmin == false} para \texttt{rol='operador'}. \\[0.3em]
El mapeo de datos debe incluir el campo rol &
  \textcolor{normalgreen}{\checkmark} Activo &
  \texttt{toMap()['rol'] == 'admin'}. \\[0.3em]
getter nombre combina todos los campos &
  \textcolor{normalgreen}{\checkmark} Nuevo &
  \texttt{nombre == 'Juan Carlos García López'}. \\[0.3em]
getter nombre omite campos opcionales vacíos &
  \textcolor{normalgreen}{\checkmark} Nuevo &
  \texttt{nombre == 'Ana Martínez'} (sin dobles espacios). \\[0.3em]
toMap incluye los cuatro campos de nombre &
  \textcolor{normalgreen}{\checkmark} Nuevo &
  \texttt{primer\_nombre}, \texttt{primer\_apellido} en el Map. \\[0.3em]
fromMap reconstruye el modelo desde nuevos campos &
  \textcolor{normalgreen}{\checkmark} Nuevo &
  \texttt{primerNombre == 'Laura'}, \texttt{primerApellido == 'Gómez'}. \\[0.3em]
fromMap usa el campo nombre legado &
  \textcolor{normalgreen}{\checkmark} Nuevo &
  Retrocompatibilidad: \texttt{primer\_nombre == ''} → usa \texttt{nombre}. \\
\end{longtable}

% ====================================================================
\section{Historial de Actualizaciones y Defectos Corregidos}
% ====================================================================

\begin{longtable}{p{2cm} p{3.5cm} p{8cm}}
\caption{Registro histórico de actualizaciones y correcciones} \label{tab:historial-updates} \\
\toprule
\textbf{Fase} & \textbf{Módulo / Área} & \textbf{Descripción del Cambio} \\
\midrule
\endfirsthead
\toprule
\textbf{Fase} & \textbf{Módulo / Área} & \textbf{Descripción del Cambio} \\
\midrule
\endhead
\bottomrule
\endlastfoot

MVP &
  Arquitectura base &
  Creación del sistema completo: simulador IoT, integración Groq API, historial SQLite,
  sistema de registro/login con validación de contraseñas. \\[0.3em]

Fase 2 (v3 BD) &
  Seguridad y SMTP &
  Migración del schema a v3 (\texttt{es\_temporal}), módulo de recuperación SMTP, algoritmo
  \texttt{generarContrasenaSegura()}, flujo \texttt{ChangePasswordScreen} con
  \texttt{PopScope(canPop:false)}. \\[0.3em]

Fase 3 / Act. \#3 &
  TDD &
  Integración de \texttt{flutter\_test}: pruebas unitarias de validación de contraseñas,
  pruebas de widget de LoginScreen, manejo de \texttt{addTearDown()} para Timers. \\[0.3em]

Act. \#4 (CRUD) &
  RBAC y Administración &
  Adición del campo \texttt{rol} en BD (migración v2), \texttt{AdminUsersScreen} con CRUD
  completo, Sudo Mode para edición/eliminación, búsqueda dinámica en memoria,
  inyección de \texttt{SECRET\_ADMIN\_CODE} vía \texttt{flutter\_dotenv}. \\[0.3em]

\textbf{Act. \#5} &
  \textbf{Registro y Flujos} &
  \textbf{(1) Separación de \texttt{nombre} en 4 campos; (2) validación de correo
  con regex; (3) navegación post-registro limpia al login; (4) cierre de sesión
  y redirección al login tras cambio de contraseña temporal; (5) migración BD v3
  automática; (6) actualización CRUD admin con 4 campos de nombre; (7) 5 nuevos
  tests unitarios.} \\

\end{longtable}

% ====================================================================
\section{Consideraciones de Seguridad y Trabajo Futuro}
% ====================================================================

\subsection{Estado Actual de Seguridad}

\begin{longtable}{p{4cm} p{4cm} p{5.5cm}}
\caption{Análisis del estado de seguridad actual} \label{tab:seguridad} \\
\toprule
\textbf{Área} & \textbf{Estado Actual} & \textbf{Mejora Recomendada} \\
\midrule
\endfirsthead
\toprule
\textbf{Área} & \textbf{Estado Actual} & \textbf{Mejora Recomendada} \\
\midrule
\endhead
\bottomrule
\endlastfoot
Almacenamiento de contraseñas &
  Texto plano en SQLite local &
  Implementar hashing con bcrypt o Argon2. \\[0.3em]
Validación de correo &
  Regex en cliente (UI) &
  Añadir restricción CHECK en SQLite o validación en backend. \\[0.3em]
Autenticación &
  Comparación directa de strings &
  Migrar a tokens JWT o sessions con expiración. \\[0.3em]
Sesión post-cambio de clave &
  \textcolor{normalgreen}{Corregido en \#5}: sesión limpiada, redirect al login &
  Considerar tiempo de expiración de sesión. \\[0.3em]
Navegación post-registro &
  \textcolor{normalgreen}{Corregido en \#5}: stack limpiado, siempre al login &
  Sin cambios adicionales requeridos. \\[0.3em]
Sudo Mode CRUD &
  Comparación de contraseña en texto plano &
  Migrar a validación con hash cuando se implemente bcrypt. \\[0.3em]
Variables de entorno &
  \texttt{flutter\_dotenv} (\texttt{.env} local) &
  En producción, usar gestión de secretos (ej. AWS Secrets Manager). \\
\end{longtable}

\subsection{Roadmap Técnico Sugerido}

\begin{enumerate}[leftmargin=2em]
  \item \textbf{Hashing de contraseñas:} Implementar \texttt{bcrypt} o \texttt{Argon2}
        antes de cualquier despliegue en producción.
  \item \textbf{Backend remoto:} Migrar de SQLite local a un servidor REST (Node.js, Django)
        con PostgreSQL para soporte multiusuario real.
  \item \textbf{Hardware físico:} Integrar el sensor YF-S201 real mediante comunicación
        BLE (Bluetooth Low Energy) o Wi-Fi para reemplazar el gemelo digital.
  \item \textbf{Notificaciones push:} Implementar alertas en tiempo real (Firebase Cloud
        Messaging) cuando la IA detecte una fuga crítica.
  \item \textbf{Dashboard analítico:} Añadir gráficas históricas (con \texttt{fl\_chart})
        para visualizar tendencias de caudal a lo largo del tiempo.
\end{enumerate}

\end{document}
